\documentclass{ibetest}
\begin{document}
\maketest{Progress Test 2.C --- Elec-S1}
         {2.C \,$\cdot$\, Pouillet's law}
         {7 minutes}{14}

% ---------------------------------------------------------------------------
\exercise{Unit conversion}{1 mark}
Tables often give resistivities in $\mu\Omega\cdot$cm. Convert
$\rho = \qty{1.7}{\micro\ohm\centi\metre}$ into $\Omega\cdot$m. Give the exact result.
\answer{1}{2.2cm}{%
  $\rho = 1.7\times10^{-6}\ \Omega\times10^{-2}\ \mathrm{m}$ \qquad
  $\boxed{\rho = 1.7\times10^{-8}\ \Omega\,\mathrm{m}}$ \quad (the resistivity of copper)}
\begin{scheme}
\pt{1}{1 pt}{$1.7\times10^{-8}\ \Omega\,\mathrm{m}$.}
\end{scheme}

% ---------------------------------------------------------------------------
\exercise{Pouillet's law}{5 marks}
State Pouillet's law giving the resistance $R$ of a conducting wire of length $L$ and
cross-sectional area $S$, made of a material of resistivity $\rho$. Give the SI unit of each
of the four quantities. Define the conductivity $\sigma$ and give its unit.
\answer{2,2,1}{5.8cm}{%
  $\boxed{R = \rho\,\dfrac{L}{S}}$ \qquad a longer wire resists more, a thicker wire
  conducts more easily.\\[5pt]
  Units: $R$ in ohms ($\Omega$), \ $\rho$ in $\Omega\,\mathrm{m}$, \ $L$ in metres (m), \
  $S$ in square metres ($\mathrm{m^2}$).\\
  \hspace*{5mm}(check: $\Omega\,\mathrm{m}\times\mathrm{m}/\mathrm{m^2} = \Omega$)\\[5pt]
  Conductivity: $\boxed{\sigma = \dfrac{1}{\rho}}$, in siemens per metre
  ($\mathrm{S\,m^{-1}}$).}
\begin{scheme}
\pt{1}{2 pts}{$R = \rho L/S$ (0 if $L$ and $S$ are swapped).}
\pt{2}{2 pts}{The four units (1~pt lost per mistake).}
\pt{3}{1 pt}{$\sigma = 1/\rho$ in $\mathrm{S\,m^{-1}}$.}
\end{scheme}
\begin{tip}
Checking the dimensions ($\Omega\,\mathrm{m}\times\mathrm{m}/\mathrm{m^2} = \Omega$)
catches a swapped $L$ and $S$ in two seconds.
\end{tip}

% ---------------------------------------------------------------------------
\exercise{Resistivity and geometry}{3 marks}
\question{A wire is replaced by a wire of the same material, twice as long and with twice the
diameter. Its resistance is:}
\opts{0}{multiplied by 2}{0}{unchanged}{1}{divided by 2}{0}{divided by 4}
\why{$L \to 2L$, and doubling the diameter multiplies $S = \pi d^2/4$ by 4, so
  $R \to R\times 2/4 = R/2$.}
\question{Among these materials, the best conductor is:}
\opts{0}{iron}{0}{aluminium}{0}{copper}{1}{silver}
\why{it has the lowest resistivity: $\rho_{\mathrm{Ag}} = 1.6\times10^{-8}\ \Omega\,
  \mathrm{m}$ (Cu: $1.7$, Al: $2.8$, Fe: $10\times10^{-8}$).}
\question{When a metallic wire heats up, its resistance:}
\opts{1}{increases}{0}{decreases}{0}{stays the same}{0}{vanishes}
\why{for most metals the resistivity increases with temperature: a heated metal conducts
  slightly less well.}

% ---------------------------------------------------------------------------
\exercise{Resistance of a copper wire}{5 marks}
A copper wire of length $L$ and cross-sectional area $S$ carries a current $i$. Determine
its resistance $R$, then the power $p$ it dissipates by the Joule effect.
\data{$L = \qty{20}{m}$ \hspace{9mm} $S = \qty{0.50}{\milli\metre\squared}$ \hspace{9mm}
      $\rho = \qty{1.7e-8}{\ohm\metre}$ \hspace{9mm} $i = \qty{5.0}{A}$}
\answer{1,1,1,1,1}{6.4cm}{%
  $\boxed{R = \rho\,\dfrac{L}{S}}$ \qquad with \
  $S = 0.50\ \mathrm{mm^2} = 0.50\times(10^{-3}\ \mathrm{m})^2 = 5.0\times10^{-7}\
   \mathrm{m^2}$\\[4pt]
  $R = \dfrac{1.7\times10^{-8}\ \Omega\,\mathrm{m}\times20\ \mathrm{m}}
            {5.0\times10^{-7}\ \mathrm{m^2}}
     = \dfrac{1.7\times20}{5.0}\times10^{-8+7}\ \Omega = 6.8\times10^{-1}\ \Omega$
  \qquad $\boxed{R = 0.68\ \Omega}$\\[4pt]
  $\boxed{p = R\,i^2} = 0.68\ \Omega\times(5.0\ \mathrm{A})^2 = 0.68\times25\ \mathrm{W}$
  \qquad $\boxed{p = 17\ \mathrm{W}}$}
\begin{scheme}
\pt{1}{1 pt}{Literal expression $R = \rho L/S$.}
\pt{2}{1 pt}{$S = 5.0\times10^{-7}\ \mathrm{m^2}$ ($1\ \mathrm{mm^2} = 10^{-6}\
  \mathrm{m^2}$).}
\pt{3}{1 pt}{$R = 0.68\ \Omega$.}
\pt{4}{1 pt}{Literal expression $p = Ri^2$.}
\pt{5}{1 pt}{$p = 17\ \mathrm{W}$, with its unit.}
\end{scheme}
\begin{tip}
$\mathrm{mm^2} \to \mathrm{m^2}$: square the conversion factor, $(10^{-3})^2 = 10^{-6}$. You
got this right in Test~1; keep writing it out explicitly.
\end{tip}

\finish
\end{document}
